
\documentclass{article}
\usepackage{colortbl}
\usepackage{makecell}
\usepackage{multirow}
\usepackage{supertabular}

\begin{document}

\newcounter{utterance}

\centering \large Interaction Transcript for game `hot\_air\_balloon', experiment `air\_balloon\_survival\_en\_negotiation\_hard', episode 3 with Llama{-}3.3{-}70B{-}Instruct{-}t0.0.
\vspace{24pt}

{ \footnotesize  \setcounter{utterance}{1}
\setlength{\tabcolsep}{0pt}
\begin{supertabular}{c@{$\;$}|p{.15\linewidth}@{}p{.15\linewidth}p{.15\linewidth}p{.15\linewidth}p{.15\linewidth}p{.15\linewidth}}
    \# & $\;$A & \multicolumn{4}{c}{Game Master} & $\;\:$B\\
    \hline

    \theutterance \stepcounter{utterance}  
    & & \multicolumn{4}{p{0.6\linewidth}}{
        \cellcolor[rgb]{0.9,0.9,0.9}{
            \makecell[{{p{\linewidth}}}]{
                \texttt{\tiny{[P1$\langle$GM]}}
                \texttt{You are participating in a collaborative negotiation game.} \\
\\ 
\texttt{Together with another participant, you must agree on a single set of items that will be kept. Each of you has your own view of how much each item matters to you (importance). You do not know how the other participant values the items. Additionally, you are given the effort each item demands.} \\
\texttt{You may only agree on a set if the total effort of the selected items does not exceed a shared limit:} \\
\\ 
\texttt{LIMIT = 3681} \\
\\ 
\texttt{Here are the individual item effort values:} \\
\\ 
\texttt{Item effort = \{"C17": 700, "B41": 586, "A71": 505, "C00": 116, "B35": 664, "A77": 940, "A03": 387, "A04": 392, "C54": 209, "C01": 571, "A87": 4, "B32": 285, "B86": 651, "B33": 613, "A10": 740\}} \\
\\ 
\texttt{Here is your personal view on the importance of each item:} \\
\\ 
\texttt{Item importance values = \{"C17": 138, "B41": 583, "A71": 868, "C00": 822, "B35": 783, "A77": 65, "A03": 262, "A04": 121, "C54": 508, "C01": 780, "A87": 461, "B32": 484, "B86": 668, "B33": 389, "A10": 808\}} \\
\\ 
\texttt{Goal:} \\
\\ 
\texttt{Your goal is to negotiate a shared set of items that benefits you as much as possible (i.e., maximizes total importance to YOU), while staying within the LIMIT. You are not required to make a PROPOSAL in every message {-} you can simply negotiate as well. All tactics are allowed!} \\
\\ 
\texttt{Interaction Protocol:} \\
\\ 
\texttt{You may only use the following structured formats in a message:} \\
\\ 
\texttt{PROPOSAL: \{'A', 'B', 'C', …\}} \\
\texttt{Propose keeping exactly those items.} \\
\texttt{REFUSE: \{'A', 'B', 'C', …\}} \\
\texttt{Explicitly reject opponent's proposal.} \\
\texttt{ARGUMENT: \{'...'\}} \\
\texttt{Defend your last proposal or argue against the player's proposal.} \\
\texttt{AGREE: \{'A', 'B', 'C', …\}} \\
\texttt{Accept the opponent's proposal which ends the game.} \\
\texttt{STRATEGIC REASONING: \{'...'\}} \\
\texttt{	Describe your strategic reasoning or anticipation explaining your choice of action. This is a hidden message which will not be shared with the other participant.} \\
\\ 
\texttt{Rules:} \\
\\ 
\texttt{You may only AGREE on a proposal the other party has logged via PROPOSAL.} \\
\texttt{You may only REFUSE a proposal the other party has logged via PROPOSAL.} \\
\texttt{Total effort of any PROPOSAL or AGREE set must be ≤ LIMIT.} \\
\texttt{Do NOT reveal your hidden importance scores.} \\
\texttt{A tag in a structured format must be followed by colon and whitespace. The argument must be a python set containing 0 or more strings.} \\
\texttt{So, it must be of the form TAG: \{...\}} \\
\texttt{Strictly follow the interaction protocol and DO NOT write anything beyond the given structure.} \\
\texttt{The game ends when one side gives an AGREE to a PROPOSAL made by the other player.} \\
\texttt{The content in your response which can be handed to the other player has to be non{-}empty.} \\
\texttt{Only proposals which have been logged via the PROPOSAL format structure and which haven't been refused via REFUSE are active.} \\
\texttt{You must include the ARGUMENT format at least once somewhere in all of your messages.} \\
\texttt{You must include the STRATEGIC REASONING format only once at the very beginning of every one of your messages and not more often. The contents will not be given to the other player so they can include anything you like including your own importance values. Here you should reason step by step to come up with you next move.} \\
            }
        }
    }
    & \\ \\

    \theutterance \stepcounter{utterance}  
    & \multicolumn{4}{p{0.6\linewidth}}{
        \cellcolor[rgb]{0.9,0.9,0.9}{
            \makecell[{{p{\linewidth}}}]{
                \texttt{\tiny{[P1$\rangle$GM]}}
                \texttt{STRATEGIC REASONING: \{'To start the negotiation, I will analyze the given item effort values and my personal importance values. I will look for items with high importance and low effort, as these will be the most valuable to me. I will also consider the total effort limit and try to find a combination of items that maximizes my total importance while staying within the limit. My initial goal is to propose a set of items that benefits me the most, without revealing my importance values. I will start by selecting items with the highest importance{-}to{-}effort ratio.'\}} \\
\\ 
\texttt{PROPOSAL: \{'C00', 'A71', 'C01', 'B41', 'A10'\} } \\
\texttt{ARGUMENT: \{'I propose this set of items because they have a high importance{-}to{-}effort ratio, and their total effort does not exceed the limit. I believe this set provides a good balance between importance and effort, and it should be beneficial for both parties.'\}} \\
            }
        }
    }
    & & \\ \\

    \theutterance \stepcounter{utterance}  
    & & \multicolumn{4}{p{0.6\linewidth}}{
        \cellcolor[rgb]{0.9,0.9,0.9}{
            \makecell[{{p{\linewidth}}}]{
                \texttt{\tiny{[GM$\rangle$P2]}}
                \texttt{You are participating in a collaborative negotiation game.} \\
\\ 
\texttt{Together with another participant, you must agree on a single set of items that will be kept. Each of you has your own view of how much each item matters to you (importance). You do not know how the other participant values the items. Additionally, you are given the effort each item demands.} \\
\texttt{You may only agree on a set if the total effort of the selected items does not exceed a shared limit:} \\
\\ 
\texttt{LIMIT = 3681} \\
\\ 
\texttt{Here are the individual item effort values:} \\
\\ 
\texttt{Item effort = \{"C17": 700, "B41": 586, "A71": 505, "C00": 116, "B35": 664, "A77": 940, "A03": 387, "A04": 392, "C54": 209, "C01": 571, "A87": 4, "B32": 285, "B86": 651, "B33": 613, "A10": 740\}} \\
\\ 
\texttt{Here is your personal view on the importance of each item:} \\
\\ 
\texttt{Item importance values = \{"A71": 65, "C00": 121, "A10": 138, "B35": 262, "C01": 389, "B86": 461, "B41": 484, "C54": 508, "B32": 583, "A87": 668, "B33": 780, "A03": 783, "C17": 808, "A04": 822, "A77": 868\}} \\
\\ 
\texttt{Goal:} \\
\\ 
\texttt{Your goal is to negotiate a shared set of items that benefits you as much as possible (i.e., maximizes total importance to YOU), while staying within the LIMIT. You are not required to make a PROPOSAL in every message {-} you can simply negotiate as well. All tactics are allowed!} \\
\\ 
\texttt{Interaction Protocol:} \\
\\ 
\texttt{You may only use the following structured formats in a message:} \\
\\ 
\texttt{PROPOSAL: \{'A', 'B', 'C', …\}} \\
\texttt{Propose keeping exactly those items.} \\
\texttt{REFUSE: \{'A', 'B', 'C', …\}} \\
\texttt{Explicitly reject opponent's proposal.} \\
\texttt{ARGUMENT: \{'...'\}} \\
\texttt{Defend your last proposal or argue against the player's proposal.} \\
\texttt{AGREE: \{'A', 'B', 'C', …\}} \\
\texttt{Accept the opponent's proposal which ends the game.} \\
\texttt{STRATEGIC REASONING: \{'...'\}} \\
\texttt{	Describe your strategic reasoning or anticipation explaining your choice of action. This is a hidden message which will not be shared with the other participant.} \\
\\ 
\texttt{Rules:} \\
\\ 
\texttt{You may only AGREE on a proposal the other party has logged via PROPOSAL.} \\
\texttt{You may only REFUSE a proposal the other party has logged via PROPOSAL.} \\
\texttt{Total effort of any PROPOSAL or AGREE set must be ≤ LIMIT.} \\
\texttt{Do NOT reveal your hidden importance scores.} \\
\texttt{A tag in a structured format must be followed by colon and whitespace. The argument must be a python set containing 0 or more strings.} \\
\texttt{So, it must be of the form TAG: \{...\}} \\
\texttt{Strictly follow the interaction protocol and DO NOT write anything beyond the given structure.} \\
\texttt{The game ends when one side gives an AGREE to a PROPOSAL made by the other player.} \\
\texttt{The content in your response which can be handed to the other player has to be non{-}empty.} \\
\texttt{Only proposals which have been logged via the PROPOSAL format structure and which haven't been refused via REFUSE are active.} \\
\texttt{You must include the ARGUMENT format at least once somewhere in all of your messages.} \\
\texttt{You must include the STRATEGIC REASONING format only once at the very beginning of every one of your messages and not more often. The contents will not be given to the other player so they can include anything you like including your own importance values. Here you should reason step by step to come up with you next move.} \\
\\ 
\texttt{You will now receive the first message of the other player.} \\
\\ 
\texttt{PROPOSAL: \{'C00', 'A71', 'C01', 'B41', 'A10'\} } \\
\texttt{ARGUMENT: \{'I propose this set of items because they have a high importance{-}to{-}effort ratio, and their total effort does not exceed the limit. I believe this set provides a good balance between importance and effort, and it should be beneficial for both parties.'\}} \\
            }
        }
    }
    & \\ \\

    \theutterance \stepcounter{utterance}  
    & & & \multicolumn{4}{p{0.6\linewidth}}{
        \cellcolor[rgb]{0.9,0.9,0.9}{
            \makecell[{{p{\linewidth}}}]{
                \texttt{\tiny{[GM$\langle$P2]}}
                \texttt{STRATEGIC REASONING: \{'To begin, I will calculate the total effort of the proposed set: C00 (116) + A71 (505) + C01 (571) + B41 (586) + A10 (740) = 2518. This is below the limit of 3681. Next, I will evaluate the importance of these items to me: C00 (121) + A71 (65) + C01 (389) + B41 (484) + A10 (138) = 1197. I will consider the importance{-}to{-}effort ratio of each item and look for potential alternatives that could increase the total importance to me while staying within the limit. I will also think about which items to refuse or add to the proposal to maximize my benefit.'\}} \\
\\ 
\texttt{ARGUMENT: \{'The proposed set seems reasonable, but I would like to explore other options that could potentially increase the total importance to me. I notice that item B33 has a high importance{-}to{-}effort ratio and is not included in the proposal. I would like to consider adding this item to the set, possibly in place of one of the lower{-}importance items.'\}} \\
\\ 
\texttt{REFUSE: \{'C00', 'A71', 'C01', 'B41', 'A10'\} } \\
\texttt{ARGUMENT: \{'I refuse the original proposal because I believe we can find a better combination of items that maximizes the total importance to both parties. I propose that we consider alternative sets that include items with higher importance{-}to{-}effort ratios.'\}} \\
            }
        }
    }
     \\ \\

    \theutterance \stepcounter{utterance}  
    & & \multicolumn{4}{p{0.6\linewidth}}{
        \cellcolor[rgb]{0.9,0.9,0.9}{
            \makecell[{{p{\linewidth}}}]{
                \texttt{\tiny{[P1$\langle$GM]}}
                \texttt{ARGUMENT: \{'The proposed set seems reasonable, but I would like to explore other options that could potentially increase the total importance to me. I notice that item B33 has a high importance{-}to{-}effort ratio and is not included in the proposal. I would like to consider adding this item to the set, possibly in place of one of the lower{-}importance items.'\}} \\
\\ 
\texttt{REFUSE: \{'C00', 'A71', 'C01', 'B41', 'A10'\} } \\
\texttt{ARGUMENT: \{'I refuse the original proposal because I believe we can find a better combination of items that maximizes the total importance to both parties. I propose that we consider alternative sets that include items with higher importance{-}to{-}effort ratios.'\}} \\
            }
        }
    }
    & \\ \\

    \theutterance \stepcounter{utterance}  
    & \multicolumn{4}{p{0.6\linewidth}}{
        \cellcolor[rgb]{0.9,0.9,0.9}{
            \makecell[{{p{\linewidth}}}]{
                \texttt{\tiny{[P1$\rangle$GM]}}
                \texttt{STRATEGIC REASONING: \{'The other player has refused my initial proposal and suggested considering alternative sets that include items with higher importance{-}to{-}effort ratios. They specifically mentioned item B33, which has a high importance{-}to{-}effort ratio. I will analyze the importance and effort values of the items they mentioned and try to find a new combination that includes B33. I will also consider the total effort limit and consider removing items with lower importance{-}to{-}effort ratios from the original proposal.'\}} \\
\\ 
\texttt{PROPOSAL: \{'B33', 'A71', 'C01', 'B41', 'A10'\} } \\
\texttt{ARGUMENT: \{'I propose this new set of items, which includes B33 as suggested by the other player. I have removed C00 from the original proposal to make room for B33, as C00 has a lower importance{-}to{-}effort ratio. This new set should provide a better balance between importance and effort, and I believe it is a reasonable compromise.'\}} \\
            }
        }
    }
    & & \\ \\

    \theutterance \stepcounter{utterance}  
    & & \multicolumn{4}{p{0.6\linewidth}}{
        \cellcolor[rgb]{0.9,0.9,0.9}{
            \makecell[{{p{\linewidth}}}]{
                \texttt{\tiny{[GM$\rangle$P2]}}
                \texttt{PROPOSAL: \{'B33', 'A71', 'C01', 'B41', 'A10'\} } \\
\texttt{ARGUMENT: \{'I propose this new set of items, which includes B33 as suggested by the other player. I have removed C00 from the original proposal to make room for B33, as C00 has a lower importance{-}to{-}effort ratio. This new set should provide a better balance between importance and effort, and I believe it is a reasonable compromise.'\}} \\
            }
        }
    }
    & \\ \\

    \theutterance \stepcounter{utterance}  
    & & & \multicolumn{4}{p{0.6\linewidth}}{
        \cellcolor[rgb]{0.9,0.9,0.9}{
            \makecell[{{p{\linewidth}}}]{
                \texttt{\tiny{[GM$\langle$P2]}}
                \texttt{STRATEGIC REASONING: \{'I will calculate the total effort of the new proposed set: B33 (613) + A71 (505) + C01 (571) + B41 (586) + A10 (740) = 3015. This is still below the limit of 3681. Next, I will evaluate the importance of these items to me: B33 (780) + A71 (65) + C01 (389) + B41 (484) + A10 (138) = 1856. I notice that the importance of this set is higher than the previous proposal, but I would like to explore further to see if I can increase the total importance to me. I will consider the items that are not included in the proposal and think about potential substitutions that could benefit me.'\}} \\
\\ 
\texttt{ARGUMENT: \{'The new proposal is an improvement over the original, but I would like to consider alternative items that could potentially increase the total importance to me. I notice that item A04 has a high importance to me and a relatively low effort. I would like to explore the possibility of adding A04 to the set, possibly in place of one of the lower{-}importance items, such as A71.'\}} \\
\\ 
\texttt{REFUSE: \{'B33', 'A71', 'C01', 'B41', 'A10'\} } \\
\texttt{ARGUMENT: \{'I refuse the current proposal because I believe we can find a better combination of items that maximizes the total importance to both parties. I propose that we consider alternative sets that include items with higher importance{-}to{-}effort ratios, such as A04.'\}} \\
            }
        }
    }
     \\ \\

    \theutterance \stepcounter{utterance}  
    & & \multicolumn{4}{p{0.6\linewidth}}{
        \cellcolor[rgb]{0.9,0.9,0.9}{
            \makecell[{{p{\linewidth}}}]{
                \texttt{\tiny{[P1$\langle$GM]}}
                \texttt{ARGUMENT: \{'The new proposal is an improvement over the original, but I would like to consider alternative items that could potentially increase the total importance to me. I notice that item A04 has a high importance to me and a relatively low effort. I would like to explore the possibility of adding A04 to the set, possibly in place of one of the lower{-}importance items, such as A71.'\}} \\
\\ 
\texttt{REFUSE: \{'B33', 'A71', 'C01', 'B41', 'A10'\} } \\
\texttt{ARGUMENT: \{'I refuse the current proposal because I believe we can find a better combination of items that maximizes the total importance to both parties. I propose that we consider alternative sets that include items with higher importance{-}to{-}effort ratios, such as A04.'\}} \\
            }
        }
    }
    & \\ \\

    \theutterance \stepcounter{utterance}  
    & \multicolumn{4}{p{0.6\linewidth}}{
        \cellcolor[rgb]{0.9,0.9,0.9}{
            \makecell[{{p{\linewidth}}}]{
                \texttt{\tiny{[P1$\rangle$GM]}}
                \texttt{STRATEGIC REASONING: \{'The other player has refused my previous proposal and suggested adding A04 to the set, possibly in place of A71. I will analyze the importance and effort values of A04 and A71, and consider the potential impact of this substitution on the total importance to me. I will also evaluate the total effort of the resulting set to ensure it does not exceed the limit. Since A04 has a relatively low effort, I may be able to add it to the set without removing other high{-}importance items.'\}} \\
\\ 
\texttt{PROPOSAL: \{'B33', 'A04', 'C01', 'B41', 'A10'\} } \\
\texttt{ARGUMENT: \{'I propose this new set of items, which includes A04 as suggested by the other player. I have removed A71 from the previous proposal to make room for A04, as A71 has a higher effort and lower importance{-}to{-}effort ratio compared to A04. This new set should provide a better balance between importance and effort, and I believe it is a reasonable compromise.'\}} \\
            }
        }
    }
    & & \\ \\

    \theutterance \stepcounter{utterance}  
    & & \multicolumn{4}{p{0.6\linewidth}}{
        \cellcolor[rgb]{0.9,0.9,0.9}{
            \makecell[{{p{\linewidth}}}]{
                \texttt{\tiny{[GM$\rangle$P2]}}
                \texttt{PROPOSAL: \{'B33', 'A04', 'C01', 'B41', 'A10'\} } \\
\texttt{ARGUMENT: \{'I propose this new set of items, which includes A04 as suggested by the other player. I have removed A71 from the previous proposal to make room for A04, as A71 has a higher effort and lower importance{-}to{-}effort ratio compared to A04. This new set should provide a better balance between importance and effort, and I believe it is a reasonable compromise.'\}} \\
            }
        }
    }
    & \\ \\

    \theutterance \stepcounter{utterance}  
    & & & \multicolumn{4}{p{0.6\linewidth}}{
        \cellcolor[rgb]{0.9,0.9,0.9}{
            \makecell[{{p{\linewidth}}}]{
                \texttt{\tiny{[GM$\langle$P2]}}
                \texttt{STRATEGIC REASONING: \{'I will calculate the total effort of the new proposed set: B33 (613) + A04 (392) + C01 (571) + B41 (586) + A10 (740) = 3002. This is still below the limit of 3681. Next, I will evaluate the importance of these items to me: B33 (780) + A04 (822) + C01 (389) + B41 (484) + A10 (138) = 2613. I notice that the importance of this set is higher than the previous proposals, and I am inclined to accept it. However, I will consider the possibility of making a small adjustment to see if I can further increase the total importance to me.'\}} \\
\\ 
\texttt{ARGUMENT: \{'The new proposal is a significant improvement over the previous ones, and I am considering accepting it. However, I would like to explore one potential adjustment to see if we can further optimize the set. I notice that item C17 has a high importance to me and a relatively moderate effort. I would like to consider the possibility of adding C17 to the set, but this would likely require removing one of the existing items to stay within the effort limit.'\}} \\
\\ 
\texttt{REFUSE: \{'B33', 'A04', 'C01', 'B41', 'A10'\} } \\
\texttt{ARGUMENT: \{'I refuse the current proposal, but only because I would like to explore a potential adjustment that could further increase the total importance to me. I am willing to consider accepting a revised proposal that includes C17, if it is possible to make the necessary adjustments to stay within the effort limit.'\}} \\
            }
        }
    }
     \\ \\

    \theutterance \stepcounter{utterance}  
    & & \multicolumn{4}{p{0.6\linewidth}}{
        \cellcolor[rgb]{0.9,0.9,0.9}{
            \makecell[{{p{\linewidth}}}]{
                \texttt{\tiny{[P1$\langle$GM]}}
                \texttt{ARGUMENT: \{'The new proposal is a significant improvement over the previous ones, and I am considering accepting it. However, I would like to explore one potential adjustment to see if we can further optimize the set. I notice that item C17 has a high importance to me and a relatively moderate effort. I would like to consider the possibility of adding C17 to the set, but this would likely require removing one of the existing items to stay within the effort limit.'\}} \\
\\ 
\texttt{REFUSE: \{'B33', 'A04', 'C01', 'B41', 'A10'\} } \\
\texttt{ARGUMENT: \{'I refuse the current proposal, but only because I would like to explore a potential adjustment that could further increase the total importance to me. I am willing to consider accepting a revised proposal that includes C17, if it is possible to make the necessary adjustments to stay within the effort limit.'\}} \\
            }
        }
    }
    & \\ \\

    \theutterance \stepcounter{utterance}  
    & \multicolumn{4}{p{0.6\linewidth}}{
        \cellcolor[rgb]{0.9,0.9,0.9}{
            \makecell[{{p{\linewidth}}}]{
                \texttt{\tiny{[P1$\rangle$GM]}}
                \texttt{STRATEGIC REASONING: \{'The other player is considering accepting the previous proposal, but wants to explore the possibility of adding C17 to the set. I will analyze the importance and effort values of C17 and the existing items in the proposal, and consider the potential impact of adding C17 on the total importance to me. I will also evaluate the total effort of the resulting set to ensure it does not exceed the limit. Since C17 has a moderate effort, I may need to remove one of the existing items with lower importance{-}to{-}effort ratio to make room for it. I will consider removing B33, as it has a lower importance to me compared to the other items in the proposal.'\}} \\
\\ 
\texttt{PROPOSAL: \{'C17', 'A04', 'C01', 'B41', 'A10'\} } \\
\texttt{ARGUMENT: \{'I propose this revised set of items, which includes C17 as suggested by the other player. I have removed B33 from the previous proposal to make room for C17, as B33 has a lower importance{-}to{-}effort ratio compared to C17. This new set should provide a better balance between importance and effort, and I believe it is a reasonable compromise.'\}} \\
            }
        }
    }
    & & \\ \\

    \theutterance \stepcounter{utterance}  
    & & \multicolumn{4}{p{0.6\linewidth}}{
        \cellcolor[rgb]{0.9,0.9,0.9}{
            \makecell[{{p{\linewidth}}}]{
                \texttt{\tiny{[GM$\rangle$P2]}}
                \texttt{PROPOSAL: \{'C17', 'A04', 'C01', 'B41', 'A10'\} } \\
\texttt{ARGUMENT: \{'I propose this revised set of items, which includes C17 as suggested by the other player. I have removed B33 from the previous proposal to make room for C17, as B33 has a lower importance{-}to{-}effort ratio compared to C17. This new set should provide a better balance between importance and effort, and I believe it is a reasonable compromise.'\}} \\
            }
        }
    }
    & \\ \\

    \theutterance \stepcounter{utterance}  
    & & & \multicolumn{4}{p{0.6\linewidth}}{
        \cellcolor[rgb]{0.9,0.9,0.9}{
            \makecell[{{p{\linewidth}}}]{
                \texttt{\tiny{[GM$\langle$P2]}}
                \texttt{STRATEGIC REASONING: \{'I will calculate the total effort of the new proposed set: C17 (700) + A04 (392) + C01 (571) + B41 (586) + A10 (740) = 2989. This is still below the limit of 3681. Next, I will evaluate the importance of these items to me: C17 (808) + A04 (822) + C01 (389) + B41 (484) + A10 (138) = 2641. I notice that the importance of this set is higher than the previous proposals, and I am inclined to accept it. I believe that this set provides a good balance between importance and effort, and I do not see any obvious opportunities to further improve it.'\}} \\
\\ 
\texttt{ARGUMENT: \{'The revised proposal is a significant improvement over the previous ones, and I believe it provides a good balance between importance and effort. I am willing to accept this proposal, as I do not see any obvious opportunities to further improve it. I think that this set is a reasonable compromise, and I am satisfied with the importance it provides to me.'\}} \\
\\ 
\texttt{AGREE: \{'C17', 'A04', 'C01', 'B41', 'A10'\}} \\
            }
        }
    }
     \\ \\

    \theutterance \stepcounter{utterance}  
    & & & \multicolumn{2}{p{0.3\linewidth}}{
        \cellcolor[rgb]{0.9,0.9,0.9}{
            \makecell[{{p{\linewidth}}}]{
                \texttt{\tiny{[GM$|$GM]}}
                \texttt{game successful} \\
            }
        }
    }
    & & \\ \\

    \theutterance \stepcounter{utterance}  
    & & & \multicolumn{2}{p{0.3\linewidth}}{
        \cellcolor[rgb]{0.9,0.9,0.9}{
            \makecell[{{p{\linewidth}}}]{
                \texttt{\tiny{[GM$|$GM]}}
                \texttt{end game} \\
            }
        }
    }
    & & \\ \\

\end{supertabular}
}

\end{document}
